% Source: https://artofproblemsolving.com/community/c3h1762011_area_in_square
% Stable identity: aops-topic-1762011-q01
Put $A=(0,0)$, $B=(1,0)$, $C=(1,1)$, $D=(0,1)$ and $h=\dfrac{\sqrt3}{2}$. The region is bounded by four arcs whose endpoints are
\[
P_1=\left(\dfrac{1}{2},1-h\right),\quad P_2=\left(h,\dfrac{1}{2}\right),\quad P_3=\left(\dfrac{1}{2},h\right),\quad P_4=\left(1-h,\dfrac{1}{2}\right);
\]
for instance $P_1$ is on the circles centred at $C$ and $D$ (it forms an equilateral triangle with $C$ and $D$), and $P_2$ is on the circles centred at $A$ and $D$. The arc $P_1P_2$ is part of the circle centred at $D$, and similarly for the other three arcs.

\begin{asy}
usepackage("amsmath");
size(6cm);
defaultpen(fontsize(10pt));
pair A=(0,0),B=(1,0),C=(1,1),D=(0,1);
real h=sqrt(3)/2;
pair P1=(0.5,1-h),P2=(h,0.5),P3=(0.5,h),P4=(1-h,0.5);
path reg=arc(D,1,-60,-30)--arc(A,1,30,60)--arc(B,1,120,150)--arc(C,1,210,240)--cycle;
fill(reg,rgb(0.8,0.88,1));
draw(arc(A,1,0,90),linewidth(0.6));draw(arc(B,1,90,180),linewidth(0.6));
draw(arc(C,1,180,270),linewidth(0.6));draw(arc(D,1,270,360),linewidth(0.6));
draw(A--B--C--D--cycle,linewidth(0.9));
label("$A$",A,SW);label("$B$",B,SE);label("$C$",C,NE);label("$D$",D,NW);
draw(P1--P2--P3--P4--cycle,dashed);
dot(P1);dot(P2);dot(P3);dot(P4);
label("$P_1$",P1,2.6*S,fontsize(8pt));label("$P_2$",P2,2.6*E,fontsize(8pt));label("$P_3$",P3,2.6*N,fontsize(8pt));label("$P_4$",P4,2.6*W,fontsize(8pt));
draw(D--P1,gray+dotted);draw(D--P2,gray+dotted);
label("\mbox{}",(min(currentpicture,user=true)+max(currentpicture,user=true))/2,black); // keeps the text after the figure black when LaTeX typesets the labels inline
\end{asy}

\textbf{The square $P_1P_2P_3P_4$.} Its side is $P_1P_2=\sqrt2\left(h-\dfrac{1}{2}\right)=\dfrac{\sqrt3-1}{\sqrt2}$, so its area is $\dfrac{(\sqrt3-1)^2}{2}=2-\sqrt3$.

\textbf{The four segments.} Seen from $D$, the directions to $P_1$ and $P_2$ make angles of $60^\circ$ and $30^\circ$ with side $DC$, so the arc $P_1P_2$ subtends $30^\circ=\dfrac{\pi}{6}$ at $D$. The circular segment between chord $P_1P_2$ and the arc has area
\[
\dfrac{1}{2}\cdot1^2\left(\dfrac{\pi}{6}-\sin\dfrac{\pi}{6}\right)=\dfrac{\pi}{12}-\dfrac{1}{4} .
\]
\textbf{Total.} The area is
\[
2-\sqrt3+4\left(\dfrac{\pi}{12}-\dfrac{1}{4}\right)=\boxed{1+\dfrac{\pi}{3}-\sqrt3}\approx0.315 .
\]
